\documentclass[11pt]{article}
\usepackage{amsmath,amssymb}
\newcommand{\Gen}[1]{\langle #1\rangle}
\title{The intersection property for the accepted marking at row 298}
\author{}
\date{}
\begin{document}
\maketitle

We retain the original degree $3510$ action and the named inner simple
group $G=\mathrm{Fi}_{22}$. The marking in question is the original free
row $298$, with owner $621$ and full witness-row digest
\texttt{ffd074b2d908ec0cdcb6633d313b665d1b43a4d71e39fcce14782356d65cb26a}.
The original model, group identification, procedural family, row and
transport identifications are the previously accepted inputs.

\paragraph{Proposition.}
The accepted marking $(a,b,c)$ at this row has the intersection property.
Consequently it is a string C-group marking of rank three.

\paragraph{Proof.}
The previous strict verification and its accepted logical generation
application establish $\Gen{a,b,c}=G$. That generation proof uses the
two supplied prime cycles and the published maximal subgroup
classification; neither is recomputed here.

The same unchanged strict worker reconstructs the complete actual endpoint
permutations and checks that they are distinct nonidentity involutions
with $ac=ca$. For this exact row it also reconstructs the entire middle
permutation $b$ and checks nonidentity, involution, and distinction from
both endpoints. These checks precede the supplied-cycle verification;
they are properties of the executed Source and its accepted case record,
not consequences inferred from a bare generation label.

All hypotheses of the existing rank-three intersection theorem in
\texttt{rank\_three\_intersections.tex} are therefore satisfied: the
ambient group is finite, nonabelian and simple; the three nonidentity
involutions generate it; and the two distinct endpoints commute.
That theorem gives, for all $I,J\subseteq\{0,1,2\}$,
\[
 \Gen{r_i:i\in I}\cap\Gen{r_i:i\in J}
   =\Gen{r_i:i\in I\cap J},
 \qquad (r_0,r_1,r_2)=(a,b,c).
\]
This is the full intersection property. The original accepted
identifications transfer it to the stated original row.
\hfill$\square$

\paragraph{Reuse and scope.}
This is a separate application of an existing theorem to an already
verified generation case. The theorem, its treatment of the intersection
of two dihedral groups, and the earlier generation proof retain their
attribution. No new group classification or new intersection theorem is
claimed. The old numerical result's zero IP credit remains unchanged;
the present proof supplies the previously omitted logical application.

The application record binds the exact theorem, prior review, actual
generation review, logical generation application, original case import,
Source and strict proof. It earns one intersection proof for an already
paid case, with no additional paid row, original manuscript, conjecture,
Lean or compression-measurement credit. Ledger registration is separate.
\end{document}
